\section{Case Studies in Modern Representation Learning}

The preceding sections developed the formal language of representations, self-supervised objectives, transfer, foundation models, and reinforcement learning.
This section studies four concrete cases through a common analytic template.
For each case we ask five questions:
\begin{enumerate}
    \item What is the training objective?
    \item What representation is learned?
    \item How is that representation transferred or adapted?
    \item What structure does the representation appear to capture?
    \item Where does the representation fail?
\end{enumerate}

\medskip

\noindent
The purpose is not to celebrate successful models in general terms.
It is to compare, in a controlled way, how different learning signals shape different kinds of reusable internal structure.

\subsection*{Case study I: BERT and masked language modeling}

A masked language model is trained by hiding selected tokens in a sequence and predicting them from context.
If a sentence is written as $x=(x_1,\dots,x_T)$ and $M\subseteq\{1,\dots,T\}$ is the set of masked positions, then a simplified masked-language-model objective is
\[
\mathcal{L}_{\mathrm{MLM}}(\theta)
=
-\sum_{i\in M} \log p_\theta(x_i\mid x_{\setminus M}).
\]
Thus the model is trained to use surrounding context to infer missing content.

\paragraph{Objective.}
The objective is predictive but not task-specific.
It does not directly optimize sentiment classification, translation, or question answering.
It instead encourages the model to estimate local and global linguistic regularities from partially observed text.

\paragraph{Learned representation.}
The learned representation is a family of contextual token embeddings
\[
h_i = \phi_\theta(x)_i,
\]
one for each token position $i$.
Unlike static word embeddings, these vectors depend on the entire sentence.
The representation of the word ``bank,'' for example, can differ substantially between a financial and a river context.
Many downstream systems also use a pooled sequence representation derived from the token states.

\paragraph{Transfer mechanism.}
Transfer usually occurs in one of two ways:
\begin{itemize}
    \item a linear probe or shallow classifier is trained on top of frozen contextual embeddings,
    \item or the full network is fine-tuned on a supervised downstream task.
\end{itemize}
The first tests how linearly accessible the task information already is; the second allows the representation itself to deform toward the new task.

\paragraph{What the representation captures.}
Empirically, masked language models learn strong syntactic and semantic structure.
Their hidden states often encode agreement, phrase boundaries, lexical disambiguation, and broad semantic similarity.
This should not be interpreted as a theorem about language understanding, but it is well supported by probing and downstream performance studies.
The important point is that contextual prediction induces a representation in which many linguistic tasks become substantially easier than in raw token space.

\paragraph{Where it fails.}
The representation also has limits.
Masked prediction does not by itself guarantee strong discourse planning, grounded reference, or robust long-horizon reasoning.
It may rely on superficial distributional cues, and it can degrade under domain shift, rare vocabulary, or tokenization mismatch.
Thus the case supports a careful conclusion:
masked prediction yields powerful reusable linguistic structure, but not a complete or theory-certified model of meaning.

\subsection*{Case study II: CLIP and contrastive image--text learning}

Contrastive image--text learning starts from paired data $(x^{\mathrm{img}},x^{\mathrm{text}})$.
An image encoder and a text encoder produce vectors
\[
z^{\mathrm{img}} = \phi_{\theta}(x^{\mathrm{img}}),
\qquad
z^{\mathrm{text}} = \psi_{\eta}(x^{\mathrm{text}}).
\]
A contrastive objective then raises the similarity of matched pairs and lowers the similarity of mismatched pairs in the same batch.
A simplified symmetric objective has the form
\[
\mathcal{L}_{\mathrm{con}}
=
-\sum_{i=1}^B
\log
\frac{\exp(\langle z^{\mathrm{img}}_i,z^{\mathrm{text}}_i\rangle/\tau)}
{\sum_{j=1}^B \exp(\langle z^{\mathrm{img}}_i,z^{\mathrm{text}}_j\rangle/\tau)}
\;+
\text{(text-to-image term)}.
\]

\paragraph{Objective.}
The objective is not to assign a fixed class label.
It is to align two modalities in a shared embedding geometry.
Learning therefore depends on correspondence across views rather than explicit task annotations.

\paragraph{Learned representation.}
The representation is a pair of embedding spaces that are trained to be mutually comparable.
Images and texts that describe the same semantic content are placed nearby.
This geometry is the central learned object.

\paragraph{Transfer mechanism.}
Transfer often occurs through prompting rather than ordinary fine-tuning.
A downstream class label such as ``cat'' is turned into a short text description, embedded by the text encoder, and compared against the image embedding.
Thus zero-shot classification becomes a nearest-neighbor or similarity-ranking problem in the shared representation space.
The same geometry also supports cross-modal retrieval.

\paragraph{What the representation captures.}
Empirically, CLIP-style models capture high-level semantic correspondences that ordinary supervised classifiers do not naturally expose.
They can connect images to flexible textual descriptions and often generalize better across open-vocabulary settings than closed-set classifiers.
This is strong evidence that contrastive cross-modal objectives can produce transferable semantic structure.

\paragraph{Where it fails.}
The learned geometry is imperfect.
CLIP-style systems can struggle with fine spatial relations, counting, OCR-heavy inputs, subtle compositional distinctions, and dataset-specific language biases.
They may also inherit social or stylistic biases from web-scale data.
Thus the case shows both the power and the fragility of alignment-based representations: they are broad, but not uniformly precise.

\subsection*{Case study III: supervised pretraining and transfer in vision}

A classical transfer-learning case begins with supervised pretraining on a large dataset such as ImageNet.
A convolutional or transformer-based vision backbone learns a representation
\[
z = \phi_\theta(x)
\]
by minimizing a labeled classification objective on the source domain.
The backbone is then reused on a smaller target dataset.

\paragraph{Objective.}
The source objective is ordinary supervised prediction.
For source examples $(x,y)$, one minimizes a loss such as
\[
\mathcal{L}_{\mathrm{sup}}(\theta,w)
=
-\log p_{\theta,w}(y\mid x).
\]
This is task-specific at training time, but the resulting representation often contains features that survive the original task boundary.

\paragraph{Learned representation.}
The learned visual representation tends to organize edges, textures, parts, object fragments, and increasingly abstract category-level information across depth.
Early layers often transfer broadly, whereas later layers may be more specialized to the source label taxonomy.

\paragraph{Transfer mechanism.}
Three standard transfer mechanisms are common:
\begin{itemize}
    \item \textbf{feature extraction:} freeze the backbone and train a new head,
    \item \textbf{linear probing:} test whether the target task is already linearly recoverable,
    \item \textbf{full fine-tuning:} update all layers on the target data.
\end{itemize}
A typical evidence-driven comparison is to evaluate these three settings on a small target dataset.
If linear probing performs well, the source representation already exposes the relevant structure.
If full fine-tuning yields a large improvement, the features are useful but not fully aligned with the target domain.

\paragraph{What the representation captures.}
This case shows that even a source objective designed for one label set can learn visual primitives useful for many other tasks.
In practice, pretrained backbones often reduce sample complexity and improve optimization on small target datasets.
This is among the clearest empirical demonstrations of representation reuse in modern machine learning.

\paragraph{Where it fails.}
Transfer is not automatic.
When the source and target domains differ sharply, for example from natural images to specialized medical images or industrial imagery, later-layer features may be poorly aligned.
A frozen backbone may miss fine texture, measurement, or domain-specific cues.
This case therefore illustrates an important general lesson:
transfer depends not only on source accuracy, but on compatibility between the source-induced representation and the target task.

\subsection*{Case study IV: deep reinforcement learning from raw observations}

In deep reinforcement learning, the system receives observations sequentially, acts, and collects rewards.
A value-based example uses a representation network together with an action-value head:
\[
z_t = \phi_\theta(x_t),
\qquad
Q_{\theta,w}(x_t,a) = g_w(z_t,a).
\]
Training then minimizes a temporal-difference loss based on bootstrapped targets such as
\[
\mathcal{L}_{\mathrm{TD}}(\theta,w)
=
\Bigl(Q_{\theta,w}(x_t,a_t) - \bigl[r_t + \gamma \max_{a'} Q_{\bar\theta,\bar w}(x_{t+1},a')\bigr]\Bigr)^2.
\]

\paragraph{Objective.}
The objective is neither supervised labeling nor ordinary self-supervised prediction.
It is control: the representation is shaped by its contribution to long-term return through value estimation or policy optimization.

\paragraph{Learned representation.}
The learned representation is a latent state useful for decision making.
From raw pixels, for example, the encoder may emphasize controllable objects, positions, velocities, or imminent hazards while suppressing irrelevant background variation.
In that sense, deep RL learns task-oriented state abstractions.

\paragraph{Transfer mechanism.}
Transfer can occur by reusing the encoder across related control tasks, initializing a new agent from a pretrained visual backbone, or combining auxiliary self-supervised objectives with policy learning.
Compared with standard transfer learning, however, reuse in RL is often less stable because the target policy changes the data distribution itself.

\paragraph{What the representation captures.}
When successful, the representation captures control-relevant geometry: which states are similar for planning, which features predict reward, and which transitions are dynamically important.
This is a distinct notion of usefulness from classification or language modeling.
A good RL representation need not be human-interpretable; it must support effective action.

\paragraph{Where it fails.}
Deep RL representations are often brittle.
They may overfit to one reward specification, fail under visual distractors or small environment changes, and learn shortcuts tied to exploration history.
Because data is generated by interaction, not passively given, representation learning is coupled to nonstationary sampling.
This makes transfer and reliability substantially harder than in standard supervised settings.

\subsection*{Comparative analysis}

These four cases support a disciplined comparison.
The choice of objective determines which invariances and distinctions are encouraged:
\begin{itemize}
    \item masked prediction emphasizes contextual reconstruction,
    \item contrastive alignment emphasizes agreement across views or modalities,
    \item supervised pretraining emphasizes source-label discrimination,
    \item reinforcement learning emphasizes control-relevant state abstraction.
\end{itemize}
The representations are therefore not interchangeable.
Each case learns a useful internal coordinate system, but the geometry of that coordinate system is shaped by the learning signal.

\begin{center}
\small
\begin{tabular}{|p{2.4cm}|p{2.8cm}|p{2.7cm}|p{2.8cm}|p{2.4cm}|}
\hline
\textbf{Case} & \textbf{Objective} & \textbf{Learned representation} & \textbf{Transfer mechanism} & \textbf{Typical failure modes} \\
\hline
BERT / MLM & Predict masked tokens from context & Contextual token and sequence embeddings & Linear probe or fine-tuning & Weak grounding, long-context limits, domain shift \\
\hline
CLIP / image--text contrast & Align matched image--text pairs, separate mismatches & Shared cross-modal embedding geometry & Prompting, zero-shot classification, retrieval & Counting, fine relations, OCR, data bias \\
\hline
Supervised vision transfer & Source-domain label prediction & Hierarchical visual features from edges to semantics & Frozen backbone, linear probe, full fine-tuning & Domain mismatch, over-specialized late features \\
\hline
Deep RL control & Maximize long-run return through bootstrapped or policy losses & Task-oriented latent state for control & Encoder reuse, auxiliary pretraining, related-task initialization & Sample inefficiency, brittleness under shift, reward dependence \\
\hline
\end{tabular}
\end{center}

\subsection*{What these cases show}

The evidence across the cases supports three chapter-level conclusions.

\paragraph{First, representation quality is objective-dependent.}
There is no single universally best representation.
The usefulness of a learned feature space depends on what information the training objective rewards the model for preserving and organizing.

\paragraph{Second, transfer is strongest when source and target structure overlap.}
A representation transfers well when the invariances and distinctions learned in pretraining remain appropriate for the downstream task.
This is why image classifiers can transfer to related visual tasks, why BERT-like encoders transfer broadly across NLP tasks, and why RL transfer is harder when rewards or transition structure change.

\paragraph{Third, failure analysis matters as much as success analysis.}
A representation should be judged not only by benchmark gains but by the structure of its errors.
Failure modes reveal what the model has and has not organized well in its internal coordinates.
This is especially important in multimodal systems and control, where apparent success can conceal brittle abstractions.

\subsection*{Notes and references}

Canonical examples include BERT for masked language modeling, CLIP for contrastive image--text learning, the large literature on ImageNet pretraining and downstream visual transfer, and deep Q-learning or actor--critic systems for learned control from high-dimensional observations.
The case-study literature is primarily empirical.
Its main contribution is comparative evidence about which objectives produce which kinds of reusable structure.
The theoretical language introduced earlier in the chapter---invariance, sufficiency, transferability, and state abstraction---provides a way to interpret that evidence more systematically.
